% =============================================================================
%  摘要与关键词
% =============================================================================
\begin{center}
{\heiti\bfseries\fontsize{14pt}{22pt}\selectfont 摘\quad 要}
\end{center}

\vspace{0.2cm}

本文以一篇基于 context-dependent DEA 的学习策略推荐论文（Expert Systems
with Applications, 2023）为对象，完成解析、复现与方法替换三项工作。首先把
原文“复合量化—聚类—类内 DEA—分层—三准则—参考路径”的证据链还原为可检验
的数学表述，指出层级不可跨群体迁移、参考路径本质是对高效同伴的模仿、障碍
度公式与表中数值不自洽等结构性局限。其次在不掌握原始问卷的条件下，依据题
项分布与锚点约束逆向重建 519 份样本，用 R 实现 PAM-Gower 聚类、类内
CCR-DEA 与分层，复现出同构的五类结构以及个体 S480 单步 26.31\%、累计
47.69\% 的提升路径。最后以随机过程替换静态评价：把五级效率层级视为状态空
间，建立带协变量的生灭过程并作极大似然估计，解析求解首达时间并用 Gillespie
模拟校验，再以 Ornstein-Uhlenbeck 扩散刻画连续效率，在学时预算下求解最优
策略。所用工具限于 Poisson 过程与一致化、Markov 链的极限定理与谱隙、鞅的停时
定理、Brown 运动的反射原理与反正弦律、Itô 公式，不涉及更新过程。
16 周末已不在第 5 层的概率 0.642、期望效率提升 25.53\% 与原文 25.38\% 相互印证；模型
还给出了原文缺失的时间尺度（中位 14.5 周，累计 43.97\% 需 53.7 周）、“增加
投入反而降低 DEA 效率”的机制性解释，以及两个只靠截面数据即可成立的结论：
从第 4 层出发先跌回第 5 层的概率为 86\%，仅靠自然波动在 16 周内触及 0.95 的
机会不足 1.5\%。

\vspace{0.4cm}

\noindent{\heiti 关键词：}随机过程；生灭过程；首达时间；数据包络分析；
context-dependent DEA；学习策略推荐
